\section{Direction Optimization That Keeps the Method Wait-Free}
\label{sec:dir-opt}

Adding bottom-up levels to the method of Section~\ref{sec:wf-topdown} raises three problems that a blocking implementation solves at its barriers. First, the direction is a decision for a whole level, which every thread must make the same way, although threads see the level at different times and can count it differently. Second, a bottom-up level works over all undiscovered vertices rather than over buckets, so it needs its own way of dividing the work and of taking over the work of a delayed thread. Third, a bottom-up level tests whether a vertex is in the frontier, and the test must stay exact for a thread that has fallen behind and runs it after the others have moved on. This section takes the three in turn. Algorithm~\ref{alg:wf-do} shows how a thread chooses the direction of each level, and Algorithm~\ref{alg:wf-bu} how it expands a bottom-up level; top-down levels are expanded as in Algorithm~\ref{alg:wf-td}, except for one change in how vertices are discovered.

\begin{algorithm}
\caption{Direction-optimizing BFS, as thread $t$ runs it. \textsc{ExpandLevel} and \textsc{AddLevel} are those of Algorithm~\ref{alg:wf-td}, with \textsc{Discover} in place of lines~\ref{ln:td-push} and~\ref{ln:td-dist}.}
\label{alg:wf-do}
\begin{algorithmic}[1]
\Procedure{BFS}{$s$}
    \State $N \gets N_0$; $P \gets$ the plan before level $0$
    \While{$N \neq \mathit{null}$}
        \State $P' \gets$ \Call{ChoosePlan}{$N$, $P$}
        \If{$P'.\mathit{direction}$ is top-down}
            \State \Call{ExpandLevel}{$N$}
        \Else
            \State \Call{ScanLevel}{$N$, $P'$, $P$}
        \EndIf
        \State $P \gets P'$; $N \gets N.\mathit{next}$
    \EndWhile
\EndProcedure
\Statex
\Procedure{ChoosePlan}{$N$, $P$}
    \If{$N.\mathit{planner}$ is empty}
        \State $N.\mathit{proposal}[t] \gets$ \Call{Rule}{$P$, $N$} \label{ln:do-propose}
        \State \Call{CAS}{$N.\mathit{planner}$, empty, $t$} \label{ln:do-claim}
    \EndIf
    \State \Return $N.\mathit{proposal}[N.\mathit{planner}]$ \label{ln:do-adopt}
\EndProcedure
\Statex
\Procedure{Discover}{$N$, $v$} \Comment{in top-down levels}
    \State append $v$ to $N.\mathit{next}.B[t]$
    \If{\Call{CAS}{$\mathit{dist}[v]$, $-1$, $N.\mathit{depth}+1$}} \label{ln:do-cas}
        \State add the degree of $v$ to $N.\mathit{next}.\mathit{edges}[t]$
    \EndIf
\EndProcedure
\end{algorithmic}
\end{algorithm}

\subsection{Agreeing on a Direction}
\label{sec:do-plan}

Before it expands a level, a thread must know the level's \emph{plan}: its direction, and for a bottom-up level the code of its frontier (Section~\ref{sec:do-codes}), together with the counts that the direction rule of Section~\ref{sec:bg-bfs} carries from one level to the next, such as $m_u$ and the size of the previous frontier. All threads must use the same plan for a level. In a top-down level, the fields $\mathit{last}$ and $\mathit{done}$ of an entry refer to positions in its bucket, while in a bottom-up level they refer to blocks of vertices, so a thread that scanned blocks while another expanded buckets would misread the other's progress, and could skip work that no thread has done. A bottom-up level must also test the frontier with the code that was written for it.

A blocking implementation lets one thread decide and makes the others wait for the decision. Instead, each node $N_d$ has a field $\mathit{planner}$, which starts out empty, and each of its entries has room for a \emph{proposal} (\textsc{ChoosePlan}). A thread that enters $N_d$ and finds $\mathit{planner}$ empty computes a plan with the direction rule (\textsc{Rule}) from what it sees: the frontier size $n_f$, which is the sum of the sizes of the buckets of $N_d$; the count $m_f$, which is the sum of the edge counts of its entries, described below; and the plan of level $d-1$. It writes the plan into its own entry (line~\ref{ln:do-propose}) and then tries to claim the level with a CAS on $\mathit{planner}$ from empty to its own number (line~\ref{ln:do-claim}). Whether its CAS succeeds or not, it then adopts the proposal of the thread that $\mathit{planner}$ names, as does every thread that finds $\mathit{planner}$ already set (line~\ref{ln:do-adopt}). A thread makes at most one CAS per level, and none waits for another: a thread claims a level only after it has written its proposal, so the proposal is there for every thread that reads the claim.

Proposals can differ, since threads read the buckets at different times and late appends change their sizes, but the CAS makes every thread adopt the same one. Which one wins does not matter for correctness, since the method finds the same distances in either direction; the rule only affects speed, and proposals made from the same buckets differ only by the few late copies. Since all threads adopt the same plan for every level, and start from the same plan before level 0, they all propose from the same plan of the level before.

The rule needs $m_f$, the number of edges of the frontier, and it gets it from the top-down level that discovered the frontier. In the method of Section~\ref{sec:wf-topdown}, every thread that finds a vertex undiscovered appends it and writes its distance, so a vertex discovered by two threads at once would have its edges counted twice. Counted that way, $m_f$ would overstate the edges of a large frontier, which can make the rule switch to bottom-up a level too early, and $m_u$, which drops by $m_f$ with every top-down level, would run down too fast. Top-down levels therefore discover a vertex with a CAS on its distance from $-1$ to $d+1$, still after appending it (\textsc{Discover}), and only the thread whose CAS succeeds adds the degree of the vertex to the edge count of its entry in $N_{d+1}$. A thread whose CAS fails has found the vertex discovered by another thread, which appended it first, so the order of Section~\ref{sec:td-order} still holds. After a bottom-up level, the rule takes $m_f = 1$, as GAPBS does.

\begin{algorithm}
\caption{A bottom-up level, as thread $t$ runs it. $P'$ is the plan of the level and $P$ that of the level before; $X_{o,k}$ is the $k$-th block of thread $o$, and $K_o$ the number of its blocks.}
\label{alg:wf-bu}
\begin{algorithmic}[1]
\Procedure{ScanLevel}{$N$, $P'$, $P$}
    \IfThen{$N.\mathit{next} = \mathit{null}$}{\Call{AddLevel}{$N$}}
    \If{$P'.\mathit{code} > 0$ \textbf{and} $P.\mathit{direction}$ is top-down}
        \State \Call{WriteCodes}{$N$, $P'.\mathit{code}$} \label{ln:bu-codes}
    \EndIf
    \For{$k \gets 0$ \textbf{to} $K_t - 1$, \textbf{while not} $N.\mathit{done}[t]$} \label{ln:bu-own}
        \IfThen{$\mathit{stamp}[t][k] \neq N.\mathit{depth}$}{\Call{ScanBlock}{$N$, $P'$, $X_{t,k}$}}
        \State $N.\mathit{last}[t] \gets k$ \label{ln:bu-last}
    \EndFor
    \State $N.\mathit{done}[t] \gets \mathit{true}$
    \For{$j \gets 1$ \textbf{to} $T-1$}
        \State $o \gets (t + j) \bmod T$
        \ForAll{blocks $k > N.\mathit{last}[o]$, in a random order}\label{ln:bu-help}
            \IfThen{$N.\mathit{done}[o]$}{\textbf{break}}
            \If{$k > N.\mathit{last}[o]$ \textbf{and} $\mathit{stamp}[o][k] \neq N.\mathit{depth}$}
                \State \Call{ScanBlock}{$N$, $P'$, $X_{o,k}$}
                \State $\mathit{stamp}[o][k] \gets N.\mathit{depth}$ \label{ln:bu-stamp}
            \EndIf
        \EndFor
        \State $N.\mathit{done}[o] \gets \mathit{true}$
    \EndFor
\EndProcedure
\Statex
\Procedure{ScanBlock}{$N$, $P'$, $X$}
    \ForAll{$v \in X$ with $\mathit{dist}[v] = -1$}
        \If{some neighbor $u$ of $v$ has $\mathit{code}[u] = P'.\mathit{code}$} \label{ln:bu-test}
            \State append $v$ to $N.\mathit{next}.B[t]$ \label{ln:bu-append}
            \State $\mathit{code}[v] \gets P'.\mathit{code} + 1$ \label{ln:bu-code}
            \State $\mathit{dist}[v] \gets N.\mathit{depth} + 1$ \label{ln:bu-dist}
        \EndIf
    \EndFor
\EndProcedure
\Statex
\Procedure{WriteCodes}{$N$, $c$}
    \For{$k \gets 0$ \textbf{to} $T-1$}
        \State $o \gets (t + k) \bmod T$
        \IfThen{$N.\mathit{codesDone}[o]$}{\textbf{continue}}
        \State $\mathit{size} \gets$ the size of $N.B[o]$
        \If{$o = t$}
            \For{$i \gets 0$ \textbf{to} $\mathit{size}-1$}
                \State $\mathit{code}[N.B[t][i]] \gets c$; $N.\mathit{codesLast}[t] \gets i$
            \EndFor
        \Else
            \For{$i \gets N.\mathit{codesLast}[o]+1$ \textbf{to} $\mathit{size}-1$}
                \State $\mathit{code}[N.B[o][i]] \gets c$
            \EndFor
        \EndIf
        \State $N.\mathit{codesDone}[o] \gets \mathit{true}$
    \EndFor
\EndProcedure
\end{algorithmic}
\end{algorithm}

\subsection{Bottom-Up Levels}
\label{sec:do-bu}

A bottom-up level divides the vertices into blocks of consecutive vertices, all of the same size, and gives block $b$ to thread $b \bmod T$, so that the blocks of every thread spread over the whole range of vertex numbers, and a level whose undiscovered vertices lie mostly in one part of the range still gives every thread a share. Blocks take the place of buckets (\textsc{ScanLevel}). A thread scans its own blocks in order, publishing in $\mathit{last}_{d,t}$ the number of the last one it has finished (line~\ref{ln:bu-last}), and then helps the other threads with the blocks after their $\mathit{last}$, in a random order of its own (line~\ref{ln:bu-help}). The flag $\mathit{done}_{d,o}$ is set once all of the blocks of thread $o$ are finished, by $o$ after its last block, or by a helper that has got through the rest; the owner and its helpers stop as soon as they see it set.

A helper cannot rely on the owner's $\mathit{last}$ alone, since a delayed owner stops moving it, and every helper of a delayed owner would then scan all of the owner's remaining blocks. A helper therefore stamps every block it finishes with the depth of the level (line~\ref{ln:bu-stamp}), and the owner and the other helpers skip a block whose stamp equals the current depth. Stamps outlive levels: a stamp left by an earlier level holds an older depth, and a thread that has fallen behind can only write the depth of its own level, which is older than the current one. Either way, a stale stamp makes a thread scan a block again, but never skip one that needs scanning.

Scanning a block means testing every vertex $v$ of the block that is still undiscovered (\textsc{ScanBlock}). If $v$ has a neighbor in the frontier, it is at distance $d+1$: the thread appends $v$ to its own bucket in $N_{d+1}$, writes the code of level $d+1$ for $v$, and only then sets $\mathit{dist}[v] = d+1$ (lines~\ref{ln:bu-append} to~\ref{ln:bu-dist}). A thread that finds $v$ discovered therefore finds it in a bucket and with its code. Threads that scan a block at the same time can discover the same vertex, but they all write the same values, so plain stores suffice. Since the discoveries of a bottom-up level go into buckets, as those of a top-down level do, the next level can go in either direction, and its plan can count its frontier.

\subsection{Frontier Codes}
\label{sec:do-codes}

The bottom-up step tests, for neighbor after neighbor, whether a vertex $u$ is in the frontier. The test could read $\mathit{dist}[u] = d$, but it is the step's most frequent memory access, and $\mathit{dist}$ has four bytes per vertex. GAPBS keeps the frontier in a bitmap instead, which it rebuilds for every bottom-up step. A bitmap reused from one level to the next would not do for our method: a thread that has fallen behind would test the frontier of level $d$ against bits that already describe a later level. We keep one byte per vertex, $\mathit{code}[v]$, and give each level that a bottom-up level tests a code of its own, $c_d$, which the plan assigns. Codes are handed out in increasing order and never reused within a search, and a vertex only ever receives the code of its own level. The test $\mathit{code}[u] = c_d$ (line~\ref{ln:bu-test}) is therefore exact at any time, for any thread, however far it has fallen behind.

A bottom-up level $d$ gets the codes of its frontier in one of two ways. If level $d-1$ was bottom-up too, its discoveries wrote $c_d$, which the plan of level $d-1$ reserved, before their distances, so every vertex of level $d$ has its code once level $d-1$ is complete. If level $d-1$ was top-down, whose appends write no codes, then every thread starts level $d$ by writing $c_d$ for the vertices in the buckets of $N_d$ (\textsc{WriteCodes}, line~\ref{ln:bu-codes}). This pass is shared the way an expansion is: a thread writes the codes of its own bucket in order, publishing how far it has got, then writes the rest of every other bucket whose codes are not done yet, and marks a bucket done once its codes are all written. Each entry has a separate pair of fields for this pass, $\mathit{codesLast}$ and $\mathit{codesDone}$, so that it does not disturb the progress of the expansion. A thread starts testing only after the pass, when every vertex of level $d$ has its code. Codes written twice or late are still right, since every vertex in the buckets of $N_d$ is in level $d$. Top-down levels write no codes, since many threads appending nearby vertices would contend for the same cache lines of $\mathit{code}$.

Codes fit in a byte, so once the plans have handed out code 255, later bottom-up levels get the code 0, which is the code of a vertex without one, and test $\mathit{dist}[u] = d$ instead; the discoveries of a level with code 255 write the code 0.

\subsection{Ending the Search}
\label{sec:do-end}

As in Section~\ref{sec:wf-topdown}, the search ends at the first level without vertices, and only a top-down level ends it. A bottom-up level always links a successor, since it cannot tell in advance whether it will discover anything, and the plan of a level whose buckets are all empty is always top-down, since every proposal for it finds $n_f = 0$. Every thread that enters that level finds its buckets empty, links no successor, and returns.
